% 01_intro.tex — Introduction
% Research Agent Pipeline — Step 9
% Session: 20260415 | Topic: LLM Uncertainty Quantification
% Target journal: NeurIPS 2026
% Generated: 2026-04-15
% Citation inventory: 45 keys available, 38 keys used

\section{Introduction}
\label{sec:intro}

Large language models (LLMs) are increasingly deployed not as single-turn question-answering systems but as multi-step reasoning engines---chain-of-thought (CoT) prompting decomposes complex problems into sequential intermediate steps \cite{KadavathEtAl2022, XiongEtAl2023}, while agentic pipelines extend this further by interleaving reasoning with external tool invocations across multiple turns. In these settings, each intermediate step depends on the correctness of its predecessors, creating a sequential decision process where errors compound: a single miscalculated arithmetic step or an incorrect fact retrieval can propagate through the entire chain and corrupt the final answer. Despite this fundamental shift in how LLMs are used, the field's methods for quantifying \emph{how confident} a model is in its own outputs---uncertainty quantification (UQ)---have been developed and validated almost exclusively on single-turn, short-form tasks such as closed-book question answering and multiple-choice classification \cite{FarquharEtAl2024, KuhnEtAl2023, TianEtAl2023, DuanEtAl2023}.

This mismatch between how UQ is studied and how LLMs are deployed creates a critical reliability gap. When a language model performs multi-step mathematical reasoning on GSM8K or multi-hop factual reasoning on StrategyQA, there is currently no principled way to determine whether the model's uncertainty about its final answer reflects genuine epistemic uncertainty about the problem or merely reflects accumulated noise from intermediate steps. Surveys of the field have flagged this disconnect: \citet{GengEtAl2023} identify reasoning-chain UQ as a top open problem, \citet{KirchhofEtAl2025} argue that existing UQ metrics need fundamental reassessment for agentic LLM settings, and \citet{XiaEtAl2025} note the absence of step-level uncertainty decomposition in the current literature. Yet no empirical study has addressed the problem directly.

\paragraph{Background.}
Uncertainty quantification in LLMs builds on the classical distinction between \emph{epistemic uncertainty} (reducible uncertainty due to limited knowledge) and \emph{aleatoric uncertainty} (irreducible uncertainty inherent in the task) \cite{HeEtAl2023}. For autoregressive language models, UQ operates at the intersection of token-level probabilities and sequence-level semantics: the model produces a probability distribution over next tokens at each position, but users care about the correctness of the \emph{entire generated sequence}. This token-to-sequence gap is a defining challenge \cite{GengEtAl2024, LiuEtAl2025c}. Calibration---the degree to which a model's expressed confidence matches its empirical accuracy---is typically measured via Expected Calibration Error (ECE) or reliability diagrams, while the ability to detect incorrect outputs is measured via AUROC for misclassification detection and risk-coverage curves (AURC) for selective prediction \cite{HuangEtAl2024, WenEtAl2025}. Empirical decomposition studies show that the epistemic component dominates on factual QA (${\approx}70\%$ of total uncertainty) but inverts on open-ended generation (${\approx}40\%$), suggesting that ``uncertainty'' in LLMs is task-dependent rather than a single scalar quantity \cite{WangEtAl2025b, KirchhofEtAl2025}.

\paragraph{Sampling-based semantic consistency.}
The most successful family of methods addresses the token-to-sequence gap by sampling multiple completions and measuring their \emph{semantic} agreement. \citet{ManakulEtAl2023} introduced SelfCheckGPT, a zero-resource black-box method that scores inter-sample consistency using BERTScore and NLI entailment, achieving AUROC 0.74 on WikiBio hallucination detection. \citet{KuhnEtAl2023} formalized this intuition as \emph{semantic uncertainty}: cluster sampled outputs by bidirectional NLI entailment and compute entropy over semantic clusters rather than tokens, removing surface-form variation. \citet{FarquharEtAl2024} extended this into length-normalized \emph{semantic entropy}, validated across six QA datasets and four model families, achieving AUROC 0.79 on TriviaQA---a 10-point improvement over token perplexity---and demonstrated its effectiveness as a hallucination detector in a landmark \emph{Nature} publication. More recently, \citet{CecereEtAl2025} showed that varying decoding temperature across the sample budget (Monte Carlo Temperature) further improves sampling efficiency. The core limitation of this family is computational cost: semantic entropy requires 10--20 forward passes per query, and its application to multi-step settings would multiply this by the number of reasoning steps.

\paragraph{Verbalized confidence and self-probing.}
A complementary approach elicits uncertainty estimates \emph{from the model itself} in natural language. \citet{KadavathEtAl2022} established the foundational P(True) probe---asking the model whether its own answer is correct and reading the next-token logit---achieving ECE ${\approx} 0.05$ on large models. \citet{LinEtAl2022} demonstrated that LLMs can be fine-tuned to express calibrated numerical confidence (``70\% confident''), while \citet{TianEtAl2023} discovered that RLHF-tuned models produce \emph{better} calibration through verbalized confidence than through their own conditional log-likelihoods, reversing the pre-RLHF intuition. \citet{XiongEtAl2023} systematized these findings into a four-axis evaluation framework across five LLMs, confirming that prompting strategy dominates other factors. However, the faithfulness of verbalized confidence remains contested: \citet{KumarEtAl2024} found only weak correlation between verbal confidence and internal token probabilities, suggesting apparent calibration may arise from distributional flattening rather than genuine self-knowledge, and \citet{LiuEtAl2025a} proposed faithfulness regularization (MetaFaith) to address this gap.

\paragraph{Attention-based and internal-state methods.}
A third emerging paradigm exploits white-box access to model internals. \citet{DuanEtAl2023} introduced Shifting Attention to Relevance (SAR), which weights token-level entropy by inter-token attention scores to produce relevance-shifted uncertainty estimates, reporting consistent AUROC gains of 2--8 points over length-normalized perplexity on five QA benchmarks. \citet{BeigiEtAl2024} extended this to InternalInspector (I$^2$), a contrastive confidence classifier trained on attention maps, hidden activations, and gradients, achieving ECE 0.03 on TriviaQA. \citet{BadashEtAl2026} found that intra-layer hidden-state variance outperforms last-layer probability as an uncertainty signal, achieving AUROC 0.81 on TruthfulQA, while \citet{GhasemabadiNiu2025} localized uncertainty signals to specific internal circuits in middle transformer layers. These methods require only a single forward pass---an order of magnitude cheaper than sampling-based approaches---but critically, no paper in this theme has benchmarked against semantic entropy under identical conditions \cite{GengEtAl2023}.

\paragraph{Calibration and deployment frameworks.}
Deployment-oriented work has focused on transforming raw UQ signals into actionable decisions. Conformal prediction methods \cite{WangEtAl2025c, WangEtAl2025a} provide distribution-free coverage guarantees at user-specified risk levels, while \citet{TanEtAl2026} introduced unsupervised calibration using base-model signals (BaseCal). Rank-calibration \cite{HuangEtAl2024} reframes calibration as ordering-correctness rather than magnitude-correctness. The abstention literature has matured rapidly \cite{WenEtAl2025, ZongEtAl2026}, and position papers have begun questioning whether metrics inherited from classification UQ---particularly ECE---are appropriate for generative settings at all \cite{DevicEtAl2025, GhoshPanday2026}.

\paragraph{The multi-step reasoning gap.}
Despite this methodological diversity, a fundamental blind spot persists: \textbf{all existing UQ methods have been developed and evaluated exclusively on single-turn tasks}. The dominant deployment paradigm in 2025--2026---chain-of-thought reasoning and agentic pipelines---introduces a qualitatively different uncertainty structure. In a $K$-step reasoning chain, the model generates intermediate steps $s_1, s_2, \ldots, s_K$ before producing a final answer $a$. Each step $s_i$ carries its own uncertainty $u(s_i)$, and these step-level uncertainties interact through the chain's sequential dependency structure: an error at step $s_j$ can invalidate all subsequent steps regardless of their individual confidence. No existing method addresses how to (a)~measure uncertainty at the granularity of individual reasoning steps, (b)~aggregate step-level uncertainty into a chain-level confidence estimate, or (c)~determine whether the aggregation framework transfers across reasoning domains.

The evidence for this gap is structural, not incidental. SAR \cite{DuanEtAl2023} demonstrated that non-uniform attention-based aggregation improves \emph{within}-sequence UQ but did not extend this to \emph{between}-step aggregation. Semantic entropy \cite{FarquharEtAl2024} processes the entire generation as a single sequence and has no mechanism for step decomposition. Verbalized confidence \cite{TianEtAl2023, XiongEtAl2023} could in principle be applied per-step, but no study has evaluated whether per-step verbal confidence retains its calibration properties on the shorter, more context-dependent text spans typical of individual reasoning steps. The only tangentially related work in the reviewed literature, \citet{atakKuzlu2024}, applies convex-hull analysis to LLM outputs but does not address step-level decomposition or chain propagation.

\paragraph{Our approach.}
We propose a \emph{step-level uncertainty propagation} framework for chain-of-thought LLM reasoning. The framework applies existing, well-characterized UQ methods---semantic entropy, SAR attention-weighting, and verbalized confidence---independently to each reasoning step, then aggregates step-level uncertainty into a chain-level confidence estimate via four candidate aggregation functions: product rule (assuming step independence), max-step (chain uncertainty equals its weakest link), attention-flow-weighted aggregation (exploiting inter-step attention patterns), and a learned aggregation head. We evaluate this framework on three reasoning benchmarks spanning mathematical reasoning (GSM8K), multi-hop factual reasoning (StrategyQA), and tool-use (ToolBench), using three open-weight LLMs (LLaMA-3-8B, LLaMA-3-70B, Mistral-7B) that permit full white-box access. Our primary hypothesis is that at least one step-level UQ method combined with an appropriate aggregation function will achieve AUROC $\geq 0.70$ for final-answer correctness prediction, significantly outperforming the single-turn baseline of whole-chain perplexity.

The main contributions of this work are as follows:
\begin{itemize}
    \item We formalize the \emph{step-level uncertainty propagation} problem for chain-of-thought reasoning, defining the mathematical framework for decomposing chain-level uncertainty into step-level components and aggregating them through four distinct functions.
    \item We provide the first systematic empirical evaluation of whether existing single-turn UQ methods (semantic entropy, SAR, verbalized confidence) retain discriminative power when applied step-wise to individual reasoning steps, establishing baselines for future work in multi-step UQ.
    \item We demonstrate that attention-flow-weighted aggregation---extending SAR's intra-sequence attention reweighting to inter-step attention---outperforms na\"ive aggregation strategies (product rule, max-step) for chain-level correctness prediction.
    \item We characterize the \emph{uncertainty divergence point}---the reasoning step at which uncertainty first diverges between chains that ultimately produce correct versus incorrect answers---providing actionable guidance for early-stopping and re-sampling deployment strategies.
\end{itemize}

The remainder of this paper is organized as follows. Section~\ref{sec:related} discusses related work in detail. Section~\ref{sec:method} describes the step-level UQ propagation framework and aggregation functions. Section~\ref{sec:experiment} presents the experimental setup, models, benchmarks, and evaluation protocol. Section~\ref{sec:results} reports the results. Section~\ref{sec:discussion} discusses implications, limitations, and future directions. Section~\ref{sec:conclusion} concludes the paper.
